\chapter{AVALIAÇÃO EXPERIMENTAL}\label{chap:avaliacao}

\section{Metodologia}

Foram executados ensaios temporários sobre o código do repositório, sem modificar
a implementação. Cada nó foi iniciado pelo sistema operacional como um processo
separado, usando o ponto de entrada \texttt{lauch\_node.py}, o catálogo original e
as portas 5001 a 5015 em \texttt{127.0.0.1}. Os comandos foram escritos na entrada
padrão dos processos em rápida sucessão e percorreram sockets TCP reais. Depois
das entregas, o comando \texttt{global} foi solicitado em todos os terminais; as
tuplas (sequência global, origem, texto) foram extraídas das saídas e comparadas.
Todos os processos encerraram normalmente.

Não há arquivos-fonte de testes no diretório \texttt{tests} da versão avaliada;
foi encontrado apenas um arquivo \texttt{.pyc} antigo. O comando
\texttt{uv run pytest} também não pôde ser usado porque \texttt{pytest} não está
declarado nas dependências. Por esse motivo, os resultados desta seção provêm do
ensaio externo descrito acima.

\section{Ordem global com 3, 8 e 15 nós}

Os três cenários foram concluídos com todos os históricos idênticos e todos os
buffers vazios, conforme a Tabela~\ref{tab:testes-ordem}. No ensaio com três nós,
cada participante emitiu uma mensagem. Nos ensaios com oito e quinze, cinco nós
emitiram simultaneamente.

\begin{table}[H]
    \centering
    \caption{Resultados da difusão e entrega ordenada}
    \label{tab:testes-ordem}
    \begin{tabular}{ccccc}
        \toprule
        \textbf{Nós} & \textbf{Envios} & \textbf{Todos receberam} &
        \textbf{Históricos idênticos} & \textbf{Buffers vazios} \\
        \midrule
        3  & 3 & Sim & Sim & Sim \\
        8  & 5 & Sim & Sim & Sim \\
        15 & 5 & Sim & Sim & Sim \\
        \bottomrule
    \end{tabular}
\end{table}

A saída resumida do cenário de 15 nós foi:

\begin{lstlisting}[numbers=none]
GLOBAL 1 | Node 1 | processo-1
GLOBAL 2 | Node 3 | processo-3
GLOBAL 3 | Node 2 | processo-2
GLOBAL 4 | Node 4 | processo-4
GLOBAL 5 | Node 5 | processo-5

comparacao dos 15 historicos: IDENTICOS
buffers restantes: [[], [], [], [], [], [], [], [], [], [], [], [], [], [], []]
\end{lstlisting}

Os pedidos dos nós 2 e 3 foram emitidos em rápida sucessão por processos
independentes, mas o líder aceitou o pedido do nó 3 primeiro. Todos os participantes
observaram a mesma inversão. Isso é uma evidência direta de que a ordem exibida foi
determinada pela sequência global e não por uma suposição local sobre a ordem de
emissão.

\section{Falha e reeleição}

Em um cenário adicional com três nós, duas mensagens foram entregues nas
sequências 1 e 2. Em seguida, o nó 1 foi encerrado e o nó 2 iniciou a eleição. Os
nós 2 e 3 convergiram para o nó 3 como novo líder, demonstrando a seleção do maior
ID ativo. Contudo, a primeira mensagem sequenciada pelo novo líder recebeu
novamente o número 1. Como os receptores sobreviventes já aguardavam a sequência
3, descartaram a mensagem como antiga. O resultado está resumido na
Tabela~\ref{tab:teste-eleicao}.

\begin{table}[H]
    \centering
    \caption{Ensaio de troca do sequenciador}
    \label{tab:teste-eleicao}
    \begin{tabular}{p{8.0cm} p{5.0cm}}
        \toprule
        \textbf{Observação} & \textbf{Resultado} \\
        \midrule
        Histórico antes da falha & $[1,2]$ em todos os nós \\
        Líder eleito entre os nós 2 e 3 & Nó 3 \\
        Primeira sequência criada pelo novo líder & 1 \\
        Próxima sequência esperada nos sobreviventes & 3 \\
        Mensagem posterior entregue & Não \\
        \bottomrule
    \end{tabular}
\end{table}

Assim, a eleição isoladamente funciona no cenário ensaiado, mas a recuperação do
serviço de ordem total não está completa. O novo coordenador precisa recuperar o
maior número estável, por exemplo consultando um quórum de históricos antes de
voltar a sequenciar.

\section{Mensagem privada}

O envio privado do nó 1 para o nó 2 abriu a conexão e retornou sucesso no emissor.
Entretanto, o vetor do receptor permaneceu $[0,0,0]$ e nenhum histórico de recepção
foi produzido. A causa, confirmada pela inspeção de \texttt{handle\_message}, é a
ausência de um ramo para \texttt{PRIVATE\_MESSAGE}. Há ainda uma mensagem de log
rotulada como privada dentro do ramo de mensagem de grupo ordenada. Portanto, o
requisito funcional de conversa privada está apenas parcialmente construído.

\section{Limitações}\label{sec:limitacoes}

Além dos resultados anteriores, a versão atual possui as seguintes limitações:

\begin{itemize}
    \item \textbf{Estado global:} não há comando, marcador nem agregação de
    snapshot; o requisito correspondente permanece pendente.
    \item \textbf{Troca de líder:} contador e histórico do sequenciador não são
    transferidos, interrompendo novas entregas após uma falha que ocorra depois de
    qualquer mensagem já entregue.
    \item \textbf{Mensagem privada:} o emissor envia o documento, mas o receptor
    não trata seu tipo.
    \item \textbf{Relógio vetorial:} em receptores remotos, uma mensagem de grupo
    ordenada atualiza o vetor durante a entrega e novamente no despachante. O líder
    local atualiza apenas uma vez, produzindo políticas de incremento diferentes.
    \item \textbf{Líder inicial:} o nó 1 é fixado como coordenador inicial, mesmo
    que o algoritmo do Valentão normalmente eleja o maior ID ativo.
    \item \textbf{TCP:} cada conexão é lida com uma única chamada de até 65.536
    bytes, sem enquadramento ou laço para remontar documentos maiores. Também não
    há autenticação nem criptografia.
    \item \textbf{Escalabilidade:} uma difusão requer uma conexão por destinatário
    e é feita sequencialmente pelo líder. O catálogo entregue termina no nó 15, e
    \texttt{start.sh} não valida pedidos acima desse limite.
    \item \textbf{Portabilidade:} a inicialização automática depende do terminal
    Konsole e de utilitários típicos de Linux. A execução manual continua possível
    em outros ambientes compatíveis com Python.
    \item \textbf{Testes:} não existe suíte reproduzível em código-fonte no estado
    atual do repositório, nem foram realizados ensaios entre máquinas físicas,
    sob perda de conexão ou com mensagens de grande volume.
\end{itemize}

\section{Síntese dos requisitos}

A Tabela~\ref{tab:requisitos} resume a cobertura verificada. ``Parcial'' indica
que existe uma parte observável do requisito, mas ao menos um cenário obrigatório
ainda falha.

\begin{table}[H]
    \centering
    \caption{Cobertura dos requisitos funcionais}
    \label{tab:requisitos}
    \begin{tabular}{c p{3.3cm} p{8.7cm}}
        \toprule
        \textbf{Req.} & \textbf{Situação} & \textbf{Evidência} \\
        \midrule
        R1 & Atendido & Difusão TCP entregue a todos em 3, 8 e 15 nós. \\
        R2 & Parcial & Chat de grupo funciona; recepção privada está ausente. \\
        R3 & Atendido & Inicialização e entrega verificadas com 3, 8 e 15 processos. \\
        R4 & Parcial & Históricos são idênticos no caminho normal; troca de líder
        perde continuidade da sequência. \\
        R5 & Parcial & Terminal oferece envio, ordens e estado, mas o unicast
        privado não é processado no destino. \\
        R6 & Não atendido & Não há captura nem exibição de estado global. \\
        R7 & Atendido & Itens documentais do roteiro são cobertos neste relatório. \\
        R8 & Parcial & O maior ID ativo é eleito, mas o novo líder não recupera a
        sequência anterior. \\
        \bottomrule
    \end{tabular}
\end{table}
